% !PRINTED 598
\weberpage{598}
The number
\weq{4}{h = (\overline{O},\, \overline{E}A)}
is called the \gesperrt{class number of the field $\Omega$ with respect to $A$}, and the system
$A\mathfrak{a}_1, A\mathfrak{a}_2, \ldots, A\mathfrak{a}_h$ the \gesperrt{ideal classes with respect to $A$}. We can prove
Theorem 2 by the general theorems on groups in \S\,100, 12.
According to this:
\[
\begin{aligned}
(\overline{O},\, \overline{E}A) &= (\overline{O},\, \overline{E}O)\,(\overline{E}O,\, \overline{E}A), \\
(\overline{E}O,\, \overline{E}A) &= (\overline{E}O,\, \overline{E}EA) = (O,\, EA), \\
(O,\, EA)\,(EA,\, A) &= (O,\, A), \\
(EA,\, A) &= (E,\, E'),
\end{aligned}
\]
if $E'$ denotes the group of numerical units contained in $A$. Consequently
\weq{5}{(\overline{O},\, \overline{E}A) = (\overline{O},\, \overline{E}O)\,\dfrac{(O,\, A)}{(E,\, E')},}
and $(\overline{O},\, \overline{E}O)$ is the number of absolute ideal classes in $\Omega$, hence
a finite number. In the quadratic field with negative discriminant
the number $(E,\, E')$ is equal to 1 if $-1$ occurs in $A$,
equal to 2 if $-1$ does not occur in $A$, and it can rise to 4 or
to 6 only in the two exceptional cases $\Delta = -4$, $\Delta = -3$.

If $\mathfrak{a}_1, \mathfrak{a}_1'$ belong to a class $\overline{A_1}$, and $\mathfrak{a}_2, \mathfrak{a}_2'$ to a class
$\overline{A_2}$, then the products $\mathfrak{a}_1 \mathfrak{a}_2$ and $\mathfrak{a}_1' \mathfrak{a}_2'$ belong to the \gesperrt{same class}, which
we call the class $\overline{A_1}\,\overline{A_2}$. The classes can therefore be composed and
form a finite \gesperrt{Abelian group}, in which the
principal class $\overline{E}A$ is the unit.

A number group $A$ shall be called a \gesperrt{congruence group} if it
consists not merely of individual numbers, but of whole number
classes with respect to a modulus $M$; that is, if $\alpha$ belongs to $A$,
then all numbers congruent to $\alpha$ with respect to the modulus $M$
shall also belong to $A$; $M$ shall be called the \gesperrt{modulus of the congruence group}.
A congruence group always has a finite index, since
the \gesperrt{number classes} with respect to the modulus $M$ are already finite in number.
Instead of $M$ one may also choose any multiple of $M$ as the modulus.
For this reason we also gain no more general concept if
we take an ideal for $M$. We may even assume $M$ to be a natural
number.

If an ideal $\mathfrak{m}$ is relatively prime to $M$, then in
every number class represented by $\alpha$ one can determine a number $\alpha_0$ from the
congruence
%CONT>>
